\chapter{Introduction}\label{ch:introduction}

\section{Motivation}

Linear-time recurrent sequence models share a common architectural
family. RWKV, Mamba-2, and Linear Attention all express their hidden
state as a bounded-dimensional recurrence, and the family has moved
into production-substrate scale on recent industrial deployments
and research efforts, including MiniMax-01~\cite{li2025minimax01},
Tencent's Hunyuan-TurboS~\cite{liu2025hunyuanturbos}, NVIDIA's
Nemotron-H series~\cite{blakeman2025nemotronh}, AI21 Labs' Jamba
hybrid Transformer-Mamba release~\cite{lieber2024jamba}, the
Technology Innovation Institute's Falcon Mamba pure-Mamba
LLM~\cite{zuo2024falconmamba}, the Evo~2 genomic foundation model
on a StripedHyena-2 linear-time substrate~\cite{brixi2026evo2},
and Apple's cross-architecture distillation recipe from
Transformer teachers to Mamba
students~\cite{moudgil2026attentiontomamba}. Underneath this
architectural variety, each backbone leaves different residual
deficits along conceptually distinct but empirically interacting
axes of expressivity, including short-range temporal mixing,
associative recall under interference, and the algebraic richness
of the recurrent transition class. The literature has been
rediscovering these same deficits for the linear-time family
across modalities (vision, audio, language, and formal
complexity), but the fixes proposed in different papers tend to
be architecture-specific and modality-bound, and have not been
compared on a single substrate at matched compute.

Each of the rediscovered deficits is, at root, a question about
\textit{expressiveness}: the capacity of a single recurrent
update or kernel to encode the state-update and data-mixing
dynamics that a task structurally requires. A short-range mixing
deficit, an associative-recall deficit, and a transition-class
limit are not three independent problems but three faces of the
same expressiveness question evaluated against three different
task structures. This thesis takes that view as its starting
point. We treat each deficit as a structural gap on a specific
axis of the diagonal-class recurrence operator, and derive a
mechanism-level extension whose form follows from the deficit's
mathematical signature. The three resulting mechanisms are
independent derivations from the deficit map developed in
Chapter~\ref{ch:theoretical}; they are not adaptations of prior
fixes. Each mechanism is then deployed on three causal backbones
at matched compute, and the transferability of its gains is read
off the cross-architecture transfer matrix.

\section{Contributions}

The thesis contributes an axis decomposition of expressivity for the
diagonal-class linear-time recurrent family, in which five
conceptually distinct but empirically interacting axes organise the
deficits each backbone leaves uncovered and provide a structural
scaffold for mechanism design. The empirical matrix exercises axes
one and two directly on its primary probes; axis three is bracketed
by the formal complexity-class bound that delimits the diagonal
class and is left to future work on the diagonal-plus-low-rank
extension. The decomposition supports the
\textit{mechanism-prior alignment} criterion, a predictive heuristic
stating that larger residual deficits along the axis a mechanism
targets tend to correspond to larger absolute gains, and that no
measured gain is expected when the axis is not exercised by the task.

The matrix that supports the criterion combines three causal
backbones (RWKV-6, Mamba-2, Linear Attention), two access modes
(causal and LION bidirectional), two parameter scales (7M and 14M),
and three probes (LibriSpeech, Common Voice English, Multi-Query
Associative Recall), populated cell by cell under a matched
50-epoch single-seed budget. The matched-budget design is what
allows the experimental chapter to attribute observed mechanism
gains to architectural cause rather than to capacity inflation, and
the single-seed convention is a deliberate scope choice under the
available compute.

The three mechanisms whose transfer the matrix measures are the
Multi-Scale Depthwise Convolution, the Damped Harmonic Oscillator,
and the Chunked Value Decorrelation, each derived from the
architectural deficit it targets rather than adapted from prior
work. The binding empirical claim is that on Multi-Query
Associative Recall at sequence length $T = 1024$, every linear-time
backbone augmented with the Multi-Scale Depthwise Convolution or
with the Chunked Value Decorrelation solves the task, while the
causal Transformer baseline fails. The result inverts the standard
sub-quadratic-versus-quadratic narrative at the sequence length
where the comparison becomes binding.

\section{Structure of the thesis}\label{sec:intro:structure}

The remainder of the thesis is organised as follows.
Chapter~\ref{ch:related} surveys the architectural and mechanism
literature on which the proposed mechanisms build, the formal
expressivity bounds that delimit the diagonal class, and the
evaluation probes used by the empirical chapter.
Chapter~\ref{ch:theoretical} establishes the theoretical
foundation, introduces the axis decomposition that frames the
research, formalises the three causal linear-time backbones
within a unified recurrence template, presents the LION parallel
form together with the decay-class taxonomy, presents the three
mathematical primitives on which the proposed mechanisms build,
and recalls the formal complexity-class bounds that delimit the
diagonal-class scope.
Chapter~\ref{ch:proposed} states the research goal, the research
questions and hypotheses, introduces the three proposed
mechanisms together with the unified LION wrapper, and describes
the diagnostic instrumentation used at the parameter level.
Chapter~\ref{ch:experiments} reports the empirical matrix, answers
the research questions across LibriSpeech, Common Voice, and
Multi-Query Associative Recall, and synthesises the
mechanism-prior alignment criterion.
Chapter~\ref{ch:conclusions} summarises the thesis's contributions
and identifies the diagonal-to-diagonal-plus-low-rank extension as
the natural next step for the axis-3 question that the present scope
does not address.
